\section{Introduction}
\label{sec:intro}

How many independent ways can board games differ from one another? This question is
not merely taxonomic --- it is foundational for any framework that claims to characterize
game design rather than rank it. If most variation in multi-axis scores is explained by
a single quality factor, then dimensional frameworks are window-dressing on a single
underlying continuum. If multiple genuinely independent dimensions exist, then games
occupy a space, and that space can be navigated.

\subsection{The General Factor Problem}

Raw principal component analysis of multi-axis game scores consistently produces a dominant
first component that we call the \emph{general quality factor}: all axes load positively
on PC1, and PC1 explains the majority of total variance~\cite{dellaLibera2026corpus}.
This is expected --- better-designed games tend to score higher on all axes simultaneously,
producing positive intercorrelations that PCA collapses into a single dimension. The
general factor is real, but it is not useful for characterizing how games differ.

\subsection{The Normalization Insight}

The central methodological contribution of this paper is the demonstration that
\emph{within-game z-score normalization} --- computing each axis score relative to
that game's own mean and standard deviation across all axes --- suppresses the general
factor and reveals the structure of residual variation. The result is a space of
design distinctiveness rather than design quality, and PCA on this normalized space
yields 4--5 genuinely independent components that correspond to theoretically coherent
design dimensions.
